\section{Reproducibility and artifact publication}\label{sec:repro}

\subsection{Two verification tiers}

The public source repository~\cite{Repository2026} separates lightweight
tracked checks from large proof replay.  Its accompanying evidence release is
designated at
\url{https://github.com/jackal092927/ramsey-bounds-reproducible/releases/tag/evidence-2026-08-30}
and is valid only if it contains exactly the 26 assets listed in
\path{artifacts/MANIFEST.tsv}.  Publication, immutability,
credential-free download, and final-tag replay are release gates recorded
separately; they are not premises of the graph theorem.

\paragraph{Tier 1: source and semantics.}
Python \(3.11.15\) and the locked environment reproduce graph verification,
unit tests, the PySAT-independent sequential-counter schemas and their frozen
block digests, and exact-seven frozen-state checks.  From the repository root,
the lightweight entry point is
\begingroup\scriptsize
\begin{verbatim}
bash scripts/bootstrap.sh
.venv/bin/python reproduce.py quick
\end{verbatim}
\endgroup
The retained maximal-union timeout endpoint is audited separately, without
recreating or promoting its deleted proof prefix, by
\begingroup\scriptsize
\begin{verbatim}
.venv/bin/python \
  routes/finite/check_r3_18_budget7_branch1_maximal_union_probe.py
\end{verbatim}
\endgroup
That checker requires exit \(124\), the exact \texttt{c UNKNOWN} result, the
fixed formula and solver hashes, and the absence of any retained file named as
a proof certificate.
The proof-critical Python packages are pinned to
\texttt{python-flint==0.9.0}, \texttt{mpmath==1.4.1},
\texttt{python-sat==1.9.dev13}, and \texttt{six==1.17.0}.  The finite route
uses PySAT; the other entries are repository-wide proof dependencies.

\paragraph{Tier 2: large proof replay.}
Once published, the release is downloaded without GitHub credentials by
\begin{verbatim}
bash scripts/download_release_artifacts.sh
\end{verbatim}
The helper first requires the complete release asset-name set to equal the
manifest exactly and then checks the content before decompression:
\begingroup\small
\begin{verbatim}
.venv/bin/python scripts/verify_artifacts.py \
  --directory /path/to/release-assets --exact
\end{verbatim}
\endgroup
The exact mode also rejects unmanifested files, so a stale or mixed download
directory cannot masquerade as the immutable release payload.
The manifest includes audited Linux x86-64 and macOS arm64
\texttt{drat-trim} binaries, their source-commit marker, and the upstream
license.  The strict branch-1 replay accepts exactly either audited binary
hash; \path{scripts/build_drat_trim.sh} remains a non-normative manual
comparison helper.  The one
integrated command
\begin{verbatim}
.venv/bin/python reproduce.py finite-heavy \
  --artifact-dir /path/to/release-assets \
  --drat-trim /absolute/path/to/audited-manifest-binary
\end{verbatim}
materializes a temporary overlay containing tracked semantic metadata and the
26 release assets.  It compares the independent sequential schemas with all
six theorem counter blocks, reconstructs the six theorem formulas and four
singleton formulas, replays all ten proofs, checks the complete common
relaxation model, and audits the exact-seven state.  A proof is accepted only when the
checker exits zero and prints an exact \texttt{s VERIFIED} status line.

\subsection{Why large proofs are release assets}

The six exact-six compressed files already occupy about \(1.26\) GB, and the
four singleton compressed proofs add several more gigabytes.  Adding these or
the smaller budget-five proofs to ordinary Git history would make cloning
needlessly expensive.
The repository therefore tracks code, matrices, cut banks, semantic records,
tests, the paper, and a content-addressed manifest.  CNF and DRAT payloads are
designated for distribution as immutable GitHub Release assets, keyed by
SHA-256 and byte count.  No DOI-bearing archival mirror is claimed.  This design keeps ordinary
clones reviewable while preserving bit-for-bit replay.

The release is complete only if it includes the three budget-five proof pairs
used by \cref{prop:budget5}, the three exact-six proof pairs used by
\cref{prop:unsat}, the four strongest singleton proof pairs, the common
CNF/model pair, and the audited checker toolchain files.  Their identities are
listed in
\cref{tab:budget5-artifacts,tab:budget6-artifacts,tab:branch1-singleton-proofs}
and in the strict manifest.  Dominated pair-core,
incomplete, corrupt, and exploratory traces are explicitly excluded.

\subsection{Reproducibility contract}

A reproduction should report five logically distinct outcomes:
\begin{enumerate}[label=(\arabic*)]
\item the seed and artifact hashes match;
\item semantic reconstruction passes;
\item each required DRAT replay is verified;
\item the common-relaxation SAT model is checked and its independent
  \(18\)-set is recovered; and
\item the exact-seven endpoint remains \textsc{unknown}, unless a new,
  separately reviewed certificate supersedes it.
\end{enumerate}
Passing only a solver run, only a hash check, or only unit tests does not by
itself reproduce the full theorem.
